\documentclass[12pt]{report}
\usepackage{graphicx}
\usepackage[letterpaper, margin={1in, 1in}]{geometry}


\begin{document}

\pagestyle{plain}

\begin{titlepage}
    \centering
    \vfill
    \includegraphics[width=4in]{img/worm_burst.jpeg}
    \vfill
    {\bfseries\Large
        Ten Page Design Document\\
        Frenzy or "Worm Game"\\
        \vskip2cm
        Bradley McFadden\\
    }
    \vfill
\end{titlepage}
~
\subsection*{Story}
Throughout history, great beasts like dinosaurs and wooly mammoths
roamed the surface of the Earth. Under the surface, even more 
intimidating creatures lurked. None were as impressive as the
sandworm, a colossal creature capable of breaching the suraface of 
the sand on a whim, and devouring large creatures.

Gaia the worm was one of these creatures, but when out hunting for
food for her offspring, but tripped up and became flash frozen in 
an icy vent underground. Millenia later, Gaia wakes in the modern
world, and ventures out to carve safe places to raise her wormlings.
Unfortunately for Gaia, apes have started walking upgright and
building machinery to harvest her underground environment. Will Gaia
be able to repopulate her race, or will she lose the fight for 
survival against the modern man?

\subsection*{Gameplay}
In \emph{Frenzy}, the player must navigate short levels, or habitats
to clear the out completely of enemies. Most of the enemies are 
miner-themed humans who will patrol and defend themselves against 
attack. Gaia is a segmented worm, and taking damage on a body part 
will split her in half. Losing enough segments or her head will cause 
her to die, and she must try again. Lucky for Gaia, she has a number
of abilities to outmanuveur and overcome the enemies infesting her
habitat, including a speed boost, digging upward or below enemies to 
use the element of surprise, and healing with a visceral bite.
Gameplay consists of navigating these levels to defeat each enemy, 
juggling ability cooldowns, using cover, and finding collectibles to
enhance her power. The game is player with a top-down camera, and levels
are divided into layers that Gaia can seamlessly switch between to take
shortcuts and gain a movement advantage over enemies.

\subsection*{Elevator pitch}
Hotline Miami meets Carrion.

\begin{figure}[h]
    \centering
    \includegraphics[width=2in]{img/worm.png}
\end{figure}

\subsection*{Game flow}
As the challenge increase and new enemies arrive, the player grows by
collecting new abilities by finding a certain unskippable collectible
and going through a short "molting" animation. Molting permanently adds
an ability to the player's save file. Players can cycle through their 
abilities using the ability bar. Up to four abilities can be mapped at
a single time. Some abilities are mutually exclusive and cannot be used
together. For example, only one bite related ability can be equipped at
one time, but players are free to switch their bite ability during a 
level. Molting will tie into the game's story by being explained as a 
part of the worm's survival skills. Intense evolutionary pressure makes
Gaia's cells search for lethal new solutions.

\begin{figure}[h]
    \centering
    \includegraphics[width=5in]{img/all.png}
\end{figure}

Other collectibles are present as well. Power-ups are included in the 
form of fossils, which can be found throughout levels. The fossils
provide themed power-ups to Gaia based on the characteristics of the 
creature that they represent. The hedgehog fossil might make it so that 
Gaia deals contact damage to enemies, or the turtle fossil might make
Gaia's segments have one hit of hyper armor before actually taking
damage. Some fossils could not be hidden, and placed in the middle of a
level to intentionally boost player powers for a particularly hard combat
section.

Another collectible is eggs, which are secrets hidden in the layer system 
or behind breakable walls of the game. Eggs unlock coats for Gaia that 
change her colour scheme, and perhaps provide minor gameplay boosts. 
For example, one egg might provide Gaia with a white themed coat that 
also makes the cooldown of the boost ability half as long.

\pagebreak
\subsection*{Character}
The player controls Gaia, a gigantic sandworm. Gaia is from the past, 
and was transported to the future by being frozen in a block of ice.
Gaia awoke today, and with no other sandworms around, it's up to her to
ensure that her species can survive in an underground world overtaken by 
man and machine. Throughout the main story, the player will exclusively
control Gaia as she seeks to reclaim her past habitats. Although, it could
make an interesting future addtion to have an inverse mode, where the player
controls a miner and defends against Gaia.

Gaia is a giant sandworm, as such she moves a large segmented body around.
Gaia has a gaping maw that can be used to attack threats and heal herself.
That's good, since enemies are fast and can quickly defeat Gaia is she is 
not careful. Gaia can die in two ways, by losing her head or falling under 
5 segments. Gaia can move up and down in the world in marked areas to attack 
from cave layers above or below her enemies. Gaia can sense movement above
her, and peek at layers above and below for an information advantage.

Gameplay is overall a mix of ducking and weaving through enemies, while
taking breaks to find collectibles.

\subsection*{Controls}
The game can be played either with a keyboard, or with a mouse and keyboard.
Mouse and keyboard steers the worm by using the mouse cursor like a carrot on
a stick, and space to move forward. Keyboard controls place both hands near the
home row. Mouse and keyboard controls could easily be adapter to also support 
a controller. The controller stick would move the cursor around for the worm to
follow. The player must always have easy access to abilities, they use them 
frequently in game, so these controls stay on A-F.

\begin{figure}[ht]
    \centering
    \includegraphics[width=3in]{img/kb_60_mouse.png}
    \includegraphics[width=3in]{img/kb.png}
\end{figure}

\pagebreak
\subsection*{Main gameplay concepts}
The player mainly engages in hit-and-run style combat, where they duck in and
away from combat while picking off enemies slowly. In this way, the game is a lot
like Hotline Miami. The player first navigates around their environment using
stealth, and when they have enough information they make a plan to take out the
10 or more enemies waiting for them. Execution is important, as they player can
be killed quickly in the presence of multiple enemies. The player's combat actions
like biting enemies makes noise, and attracts other enemies to their location,
so the player must be prepared to keep fighting once they start, or have a plan
to retreat. The player starts with two abilities to help with the duck and weave combat,
a bite that kills unarmored enemies and heals the player, and a speed boost to close
the distance or flee.

There is a system in place where the player can move underground or climb above
ground in marked areas. Most enemies cannot follow the player, so these areas
act as gates between two places that the player can use for a break or to jump
at enemies with. From the above ground or the underground, players can see 
enemy noises and look at the level geometry to plan their attacks while they
hide.

The game is split into levels that are themed for different underground environments.
There are somewhat generic caves, an ice area, volcano, underground jungle, and so
forth. Each level takes place in one of these environments, until reaching the 
end of an environment where the player will fight a boss.

\begin{figure}[h]
    \includegraphics[width=3in]{img/peek.png}
    \hfill
    \includegraphics[width=3in]{img/worm_enemy.png}
\end{figure}

The player wins by defeating all enemies in a level, and the player loses if they 
die by losing their head or enough body segments. It's not a big deal to die, levels
are meant to be about a five minute or less expericen each. The levels are also small 
enough so that the player can press a button to instantly restart without any loading
at all.

Since the game is 2D, to provide more detail to the normals they are 3D models rendered
top-down with normal maps. This allows us to light the models in an intersting way
to provide some depth and cool lighting effects for levels. Additionally, as Gaia eats
more enemies or takes more damage, her mesh changes to reflect her current health.
She becomes bruised and purple with more damage, and her mouth is stained red with blood.

\pagebreak
\subsection*{Game world}
The world of \emph{Frenzy} is based upon the rich environments envisioned by the likes of 
Deep Rock Galatic, where there is an alien presence that has invaded a variety of rich and
beautiful underground environments, to harvest it for gains. Levels are hand designed, no 
procedural generation is used. The game should reflect that the environment is starting to 
give way to the men and machinery that are harvesting it. Waters are polluted, wildlife 
looks ill. Areas with machines have some smog. Digging machines, transportation machines,
and miner habitats decorate the environment. Environments are a mix of former natural beauty
and industrialization.

There should be few different types of environments:
\begin{description}
    \item [Ice world] \hfill \\ World with vents of ice, chilling breezes that freeze enemies, and patches of ice that cause sliding. Icicles can fall from the ceiling in certain areas.
    \item [Salt flats] \hfill \\ Reddish areas with a lot of wildlife to eat the salt. Crystals are common with weird interactions.
    \item [Sandstone caves] \hfill \\ Sandstone caves with occasional storms and sand colored backgrounds. Pits can cause the worm to fall down layers. 
    \item [Mushroom world] \hfill \\ World overcomes with strange spores and fungus, with some water.
    \item [Lava world] \hfill \\ Hell world, including lava, fire traps, obsidian.
    \item [Jungle] \hfill \\ World teeming with wildlife, prey and predator.
\end{description}

Textures use shaders to warp and warble to look strange. Particle effects are also used
for snow, fire, spores etc. Some idle wildlife is around, dependent on the environment. 
Wherever people are, more machines are and more pollution is apparent. Breakable walls and
hazards are scattered appropriately. Levels do not have to be linear, and can be larger
islands from time to time. Placement of collectibles encourages the player to explore 
throughout the map to unlock everything. Secrets include hidden layer switches and breakable
decorations/walls.

\pagebreak
\subsection*{Interface}

The interface of the game is fairly standard. Starting with a skippable splash screen,
the player is brought to the title screen. From the title screen they can continue from 
the most recent level they unlocked, select a previous level, go to the credits, change 
settings, or quit out of the game. In game, a pause menu that pauses gameplay is present
that allows some settings to be changed, the user to quit back to the title screen, or quit
altogether. 

The title screen should really hint at the story of the game by showing the 
Gaia's story, from before the ice in prehistoric age, to being encased in ice, to breaking
free. The mood of the title screen should express the kind of shock of the horrible 
realization of being sort far out of time as Gaia does.

\begin{figure}[ht]
    \centering
    \includegraphics[width=5in]{img/game_flow.png}
\end{figure}

For the level select, I'm torn between a minigame style map like old Mario games, or an 
easier to navigate interface. To provide the player with some sense of what they're going 
into, I'd like to show a map of the level, as well as some kind of tally screen level. Since
the game world is going to divided into "environments" and levels within the environments, 
it would be neat to dynamically theme the level select based on the current environment.
I'd like the interface music to be mysterious.

When transitioning between game levels, the name of the current level should fade in, as 
well as the environment, perhaps with a styled border. I'd like the camera to pan and zoom 
in across the game world so the player knows about the scale of the level and has some hints
about the challenge associated with it.

\pagebreak
\subsection*{Mechanics and power-ups}
In \emph{Frenzy}, the player controls Gaia through levels with a top-down camera. The camera
follows Gaia exactly. The player can "peek" the camera ahead by about a screen to look ahead
of themselves. The player will also navigate through red/green layer exits that bring the 
layer above or below everything on their current layer. On other layers, there may be
secrets, enemies, or more collectibles to find. Layers add some space to levels and can
be used to ambush enemies or as a means of escaping enemies that can't follow.

In \emph{Frenzy}, the player has two main abilities to use to destroy their enemies, a bite 
and a speed boost. The bite ability will open the player mouth, and it auto-closes if Gaia's
head collides with an enemy. The bite will restore the player's health and add a segment
to Gaia's body. The player can track their health via a segmented health bar on the HUD. 
The bite has no cooldown. The speed boost ability grants players a burst of speed for a 
short duration. It goes on a long cooldown for several seconds, and is used to quickly 
close distance or escape. The player can keep track of what abilities are equipped with a 
abilities bar.

A noise system is present that calls enemies to the source of noises. Enemies can also be
alerted to the player's presence by Gaia moving in front of them.

Players in \emph{Frenzy} can collect eggs, fossils, and molt to gain permanent upgrades. Eggs 
are secrets that provide new "skins" for the player with minor gameplay effects. They can 
be changed out through the level select interface. Fossils grant power-ups to the player 
that are temporary, and very strong. They are sometimes hidden, and other times in plain 
sight and intended to be used as part of the challenge. Permanent power-ups are found in 
certain game levels, and can be used in a variety of ways depending on the situation. The 
player gains permanent power-ups by proceeding through an unskippable gate mechanic where an
animation plays, and the player molts. Ideally, it would be nice to see what ability or 
permanent upgrade the worm has equipped based on their appearance. The
acid spit power-up is used to melt enemies or enemy built structures,. Gaia is shown with 
green mouth patches when using the acid spit. Some upgrades only last for a specific level. 
For example, when killing a Blaster, the player can take its arm to use as a projectile 
weapon for a few shots. Other props can be interacted with sometimes for special effects. 
Engineer toolboxes can be used to breakthrough heavy doors.

Certain game levels have hazards, like timed acid traps, or gas spewing plants. Laser wires
might cut Gaia in half. Enemies can be lured into hazards to be dispatched easier. Some 
level geometry can be interacted with, such as breakable walls, or collapsible ceilings.
Certain hazards can displace the player or enemies, like a vent of gas that pushes the 
player. This could be used to the player's advantage as an additional speed boost to ambush 
enemies.

\pagebreak
\subsection*{Enemies and bosses}
The player faces off against a variety of enemies. Since the player is attacking mines 
underground, most of the enemies are mining-related, or are armed guards for the mines.

The most basic enemy the player will go up against is a standard grunt, with a knife or a 
gun. These enemies die in one hit, and are only threatening in large groups. They often 
patrol areas in a fairly regular path, but respond quickly to noises by running in their
direction.

\begin{figure}[ht]
    \centering
    \includegraphics[width=2in]{img/driller.jpeg.jpeg}
    \includegraphics[width=3in]{img/prospector.jpeg}
\end{figure}

The Driller is an enemy that Gaia will face that wields large power drills. He is heavily
armored, and cannot be felled with a bite while protected. He will hold is ground and charge
up his power drills to track the player and attempt to ram them. If he collides with a wall,
he will be vulnerable to attack by the player for a short duration. In the environment, he is 
often spotted using his drills to dig tunnels.

The Prospector is an enemy usually found wandering, looking for veins of ore. He wields a 
large pickaxe that he can use to deal large damage to Gaia, and his attacks also knock
debris from the ceiling that can also damage Gaia. He is also armored, and will take one
hit before losing his armor. When he is armor is destroyed, he will enter a new state, where 
he can throw sticks of dynamite. On death, he'll explode after a short duration.

\begin{figure}[ht]
    \centering
    \includegraphics[width=2in]{img/blaster.jpeg}
    \includegraphics[width=2in]{img/engi.jpeg}
\end{figure}

The Blaster is an enemy that mines ores with a large arm cannon. He attacks at range, and
his target is marked by a laser guide. His attacks will pierce Gaia to damage her multiple
times. If he is left alone, he can upgrade his attacks to bounce off walls. Defeating him
might allow Gaia to use his arm for herself.

The Engineer is a complicated enemy that prefers to be on the defensive. He builds breakable
walls to barricade himself in, and provides buffs to his allies, or sets up dangerous 
machines and traps. He should be targeted as soon as possible.

\pagebreak
\subsection*{Cutscenes, bonus material and comps}
There will not really be many cutscenes. A cutscene will be present at the start of any boss
fights, and at the end, as well as transitions to new worlds. Cutscenes will either be CGI 
3D renders, or puppet shows with in-game models. 

In the forseeable future, a game mode where the player plays as some of the special enemies 
could be created. Otherwise, more game worlds could be created for the player.

\begin{figure}[ht]
    \centering
    \includegraphics[width=4in]{img/dead_grunt.png}
\end{figure}

\end{document}